\hnsection{EXTENSIONS OF THE NOTION OF LIE GROUP}

Nowadays the vitality of Lie theory is manifested by the diversity of its applications (in topology, differential geometry, arithmetic, etc.) and by the creation of parallel theories where the underlying manifold structure is replaced by a similar structure ( \(p\) -adic manifold, algebraic variety, scheme, formal scheme, \(\ldots\) ). We are not concerned here with the history of these developments and we shall confine ourselves to those touched on in Chapter III: Banach Lie groups and \(p\) -adic Lie groups.

\textbf{(a) Banach Lie groups}

Here we are concerned with ``infinite-dimensional'' Lie groups. From the local point of view, a neighbourhood of 0 in Euclidean space is replaced by a neighbourhood of 0 in a Banach space. This is what G. Birkhoff did in 1936 {[}29 a{]}, thus achieving the notion of a complete normed Lie algebra and its correspondence with a ``group germ'' defined on an open set of a Banach space. About 1950, Dynkin completed these results by extending the Hausdorff formula to this case (cf.~supra).

The definitions and results of Birkhoff and Dynkin are local. Until recently, it does not seem that any one has sought to make the corresponding global theory explicit, no doubt because of the lack of applications.†

\textbf{(b) p-adic Lie groups}

Such groups were encountered for the first time in 1907 in the works of Hensel {[}19{]} on \(p\) -adic analytic functions (defined by expansions as integral series). He studied in particular the exponential and the logarithm; in spite of the a priori surprising behaviour of the series which define them (for example the exponential series does not converge everywhere), their fundamental functional properties remain valid, which provides a local isomorphism between the additive group and the multiplicative group of \(\mathbf{Q}_p\) (or, more generally, of any complete ultrametric field of characteristic zero).

A. Weil {[}33{]} and E. Lutz {[}34{]} were equally concerned with commutative (but this time non-linear) groups in their work on p-adic elliptic curves (1936). Besides arithmetic applications, a local isomorphism is constructed here between the group and the additive group, based on the integration of an invariant differential form. This method applies equally to Abelian varieties, as was noted by C. Chabauty soon afterwards, who used it without further explanation to prove a particular case of the ``Mordell conjecture'' {[}35{]}.

From then on, it was clear that the local theory of Lie groups could be applied with scarcely any change to the p-adic case. The fundamental theorems of the Lie groups--Lie algebras ``dictionary'' were established in 1942 in the thesis of R. Hooke {[}37{]}, a pupil of Chevalley; this work also contained the p-adic analogue of E. Cartan's theorem on closed subgroups of real Lie groups.

More recently, M. Lazard {[}42 b{]} developed a more precise form of the ``dictionary'' for compact analytic groups over \(\mathbf{Q}_{\mathrm{p}}\) . He showed that the existence of a \(p\) -adic analytic structure on a compact group G is closely related to that of certain filtrations on G and gives various applications of this (for example to the cohomology of G). One of Lazard's tools is an improvement of Dynkin's results on the convergence of the \(p\) -adic Hausdorff series {[}42 a{]}.

